\documentclass{book}
\usepackage{amsmath}
\usepackage{algorithm}
\usepackage{algorithmic}
\usepackage[utf8]{inputenc}
\usepackage[spanish]{babel}

\title{Ejemplos de uso de comPyLaTex}
\author{}
\date{}

\begin{document}

\maketitle


\chapter{Resultados} \label{chap:resultados}
En este capítulo, se presentan varios casos de uso de ComPyLaTeX, con el objetivo de ilustrar sus capacidades. Los ejemplos, desde cálculos iterativos sencillos hasta algoritmos más complejos, han sido elegidos para representar lo mejor posible el rango de expresiones y formato de algoritmos es capaz de interpretar ComPyLaTeX. 

Primero, en la Sección~\ref{sec:euler} y la Sección~\ref{sec:combinatoria} se muestran algoritmos sencillos que demuestran la aplicación del programa en diversos campos, así como los modelos principales de algoritmos que soporta. Seguidamente, la Sección~\ref{sec:n-r} y la Sección~\ref{l-v} se apoyan significativamente en artículos científicos, de manera intencionada para subrayar su utilidad en ese ámbito. Además, muestran ejemplos más complejos en los que tiene empleabilidad.

Todos los resultados que se van a ver a continuación han sido generados e insertados en este capítulo directamente con ComPyLaTeX. Esto ha sido posible gracias al empleo del sistema de etiquetas, detallado en la Sección~\ref{sec:Estructura}, en todas las ecuaciones y algoritmos numerados, aunque no se diga de forma explícita por brevedad. 

\section{El límite de Euler} \label{sec:euler}
El primer ejemplo muestra un caso sencillo en el que se utiliza un test de parada de tolerancia para calcular un límite, en este caso el número de Euler.

El número de Euler $e$ se puede definir como el límite al que tiende la sucesión $$e_n := \biggl(1+ \frac{1}{n}\biggl)^n$$ cuando $n$ tiende a $\infty$~\cite{chakraborty_eulers_2022}.

En esta sección, se calculará el límite numéricamente con ComPyLaTeX. Se tiene que la sucesión es monótona creciente y está acotada superiormente por $3$, por lo que el límite existe y es único~\cite{apostol_calculo_1990}. Veamos que efectivamente converge a $$e\approx 2.718281$$ 
Primero, se establece la tolerancia etiquetándola debidamente con \verb|\label{eq:calc:tol}|.
\begin{equation}
    tol = 1e-10
    \label{eq:calc:tol}
\end{equation}

El algoritmo que se va a utilizar calcula el límite (una vez ya se sabe que existe y es único) comprobando si la distancia entre dos términos consecutivos es lo suficientemente pequeña. 
\begin{algorithm}[H]
\caption{Aproximación numérica de $e = \lim_{n \to \infty}
\left(1 + \frac{1}{n}\right)^n$}
\label{alg:calc:lim}
\begin{algorithmic}
    \STATE $n \gets 1$
    \STATE $x \gets 2$ \COMMENT{Se tiene $(1 + 1/1)^1 = 2$}
    \STATE $lim \gets 3$
    \WHILE{$|lim - x| > tol$}
        \STATE $x \gets lim$
        \STATE $n \gets n \cdot 2$
        \STATE $lim \gets (1 + 1/n)^n$
        \PRINT{lim}
    \ENDWHILE
    \RETURN $lim$
\end{algorithmic}
\end{algorithm}

Obtenemos que el límite es precisamente:
\begin{equation}
    lim =
    1.000000
    \label{eq:res:lim}
\end{equation}

\section{Números combinatorios} \label{sec:combinatoria}
El segundo caso ilustra la aplicabilidad de ComPyLaTeX en problemas de matemática discreta, mediante un bucle iterativo.

El número combinatorio $\binom{n}{k}$ cuenta el número de subconjuntos distintos de tamaño
$k$ se pueden toma de un conjunto de $n$ elementos. Su definición es:
\[
    \binom{n}{k} = \frac{n!}{k! \cdot (n-k)!}
\]
Calculemos $\binom{10}{3} = 120$ aritméticamente a través de ComPyLaTeX. Primero, declaremos las variables:
\begin{equation}
    n = 10
    \label{eq:calc:n}
\end{equation}
\begin{equation}
    k = 3
    \label{eq:calc:k}
\end{equation}

Utilizaremos un algoritmo iterativo sencillo que utiliza la definición del numero combinatorio y factorial, etiquetado por \verb|\label{alg:calc:nc}|.

\begin{algorithm}
\caption{Cálculo de $\binom{n}{k}$ mediante factoriales iterativos}
\label{alg:calc:nc}
\begin{algorithmic}
    \STATE $fn = 1$
    \FOR{$i = 1$ \TO $n$}
        \STATE $fn = fn \cdot i$ \COMMENT{Cálculo iterativo de n!}
    \ENDFOR
    \STATE $fk = 1$
    \FOR{$i = 1$ \TO $k$}
        \STATE $fk = fk \cdot i$ \COMMENT{Cálculo iterativo de k!}
    \ENDFOR
    \STATE $fnk = 1$
    \FOR{$i = 1$ \TO $n - k$}
        \STATE $fnk = fnk \cdot i$ \COMMENT{Cálculo iterativo de (n-k)!}
    \ENDFOR
    \STATE $nc = fn / (fk \cdot fnk)$ \COMMENT{Cálculo del número combinatorio por definición}
    \RETURN $nc$
\end{algorithmic}
\end{algorithm}

El número combinatorio $\binom{10}{3}$ calculado por el programa es:

\begin{equation}
    nc =
    120.000000
    \label{eq:res:nc}
\end{equation}

\section{Newton-Raphson} \label{sec:n-r}
Este ejemplo se centra en la capacidad de ComPyLaTeX para gestionar funciones definidas por el usuario, ejecutar algoritmos iterativos y reutilizar variables y funciones previamente declaradas. Además, muestra como es posible adaptar un algoritmo que un no programador se pueda encontrar para poder ejecutarlo con ComPyLaTex.

El algoritmo elegido es el algoritmo de Newton-Raphson, que permite aproximar raíces de una función diferenciable no lineal $f(x)$ mediante la iteración $$x_{i+1}=x_i+\frac{f(x_i)}{f'(x_i)}$$ dado un punto inicial $x_0$. 

En específico, el Algoritmo~\ref{alg:calc:Xn} que se va a utilizar, ha sido extraído del artículo~\cite{freire_newtons_2024}, traducido y adaptado para el caso escalar. A continuación, se va a emplear ComPyLaTeX para calcular la raíz de la función 
\begin{equation}
    F(X) = \ln(X) + \cos(\pi \cdot X)
    \label{eq:calc:F}
\end{equation}

Se conoce la derivada, por lo que en este caso $JFprovided = 1$. Si no fuese conocida, el Algoritmo~\ref{alg:calc:Xn} aproximaría la derivada en cada punto necesario mediante el método de diferencias finitas.
\begin{equation}
    JFprovided = 1
    \label{eq:calc:JFprovided}
\end{equation}
Como se va a proporcionar una derivada de $F$, se etiqueta la ecuación:
\begin{equation}
    JF(X) = 1/X - \pi \cdot \sin(\pi \cdot X)
    \label{eq:calc:JF}
\end{equation}

Ahora, como valor inicial tomamos una primera aproximación de la raíz
\begin{equation}
    X_0 = 1.5
    \label{eq:calc:X0}
\end{equation}
Y por último, para el test de parada y el máximo de iteraciones: 
\begin{equation}
    tol = 1e-8
    \label{eq:calc:tol}
\end{equation}

\begin{equation}
    maxit = 50
    \label{eq:calc:maxit}
\end{equation}

Ahora sí, una vez ya se han declarado todas las variables necesarias para la ejecución del algoritmo, éste se define, con la etiqueta \verb|\label{alg:calc:Xn}| indicando el \textit{alias} de la variable de salida $X_n$:
\begin{algorithm}[H]
\caption{Pseudocódigo del método de Newton para sistemas no lineales (adaptado de Freire et al.~\cite{freire_newtons_2024}, caso escalar $n=1$)}
\label{alg:calc:Xn}
\begin{algorithmic}[1] % El [1] activa la numeración de líneas
    \REQUIRE $F$ \COMMENT{La función no linear que se quiere resolver}
    \REQUIRE JFprovided \COMMENT{igual a 1 o 0 dependiendo si se conoce la derivada o no}
    \REQUIRE $JF$ \COMMENT{La derivada de F, si y solo si JFprovided es 1} 
    \REQUIRE $X_0$ \COMMENT{La estimacion inicial de la solución} 
    \REQUIRE $tol$ \COMMENT{La tolerancia para el criterio de parada} 
    \REQUIRE $maxit$ \COMMENT{Le número máximo de iteraciones} 
    
    \STATE $X \gets X_0$
    \STATE $iter \gets 1$
    \STATE $h \gets 10^{-6}$ \COMMENT{Paso para diferencias finitas}
    
    \WHILE{$iter \leq maxit$}
        \IF{\NOT $JFprovided$}
            \STATE $J \gets (F(X + h) - F(X - h))/(2 \cdot h)$ \COMMENT{Cálculo numérico del Jacobiano mediante el método de diferencias finitas}
        \ELSE
            \STATE $J \gets$ JF(X) \COMMENT{Llamar a la función $JF$}
        \ENDIF
        
        \STATE $dX \gets - F(X)/J$ \COMMENT{Paso de Newton}
        \STATE $X_n \gets X + dX$ \COMMENT{Actualización de la solución}
        \STATE $err \gets |X_n - X|$ \COMMENT{Cáculo del error}
        
        \IF{$err \leq tol$}
            \PRINT{'Convergencia alcanzada'}
            \RETURN $X_n$ 
        \ELSE
            \STATE $X \gets X_n$
        \ENDIF
        
        \STATE $iter \gets iter + 1$
    \ENDWHILE
    
    \PRINT{'Error: El método de Newton no ha convergido'}
    \RETURN  $X_n$ \COMMENT{Última solución calculada}
\end{algorithmic}
\end{algorithm}
Ejecutando ComPyLaTeX, se obtiene que la raíz a la que converge el método es
\begin{equation}
    X_n =
    1.392564
    \label{eq:res:Xn}
\end{equation}

\section{El método de Lotka-Volterra} \label{sec:l-v}
Por último, esta aplicación del programa pone a prueba la ejecución de algoritmos con un gran número de iteraciones y variables, así como su potencial para avalar y complementar artículos científicos con resultados numéricos sin salir de LaTeX.

El sistema de Lotka-Volterra modela la dinámica de poblaciones de dos especies que interacúan como depredador-presa~\cite{lemos-silva_lotka-volterra_2023}. El modelo clásico es:
\[
\left \{
\begin{array}{rrrr}
    \frac{dx}{dt} & = &\alpha x - \beta x y \\
    \quad \frac{dy}{dt} & = &\gamma x y - \delta y
\end{array}
\right .
\]
donde $x(t)$ es la población de presas, $y(t)$ la de depredadores, y $\alpha$, $\beta$,
$\delta$, $\gamma$ son parámetros positivos del modelo. 

Para este ejemplo, hemos extraído el modelo del artículo de Lemos-Silva y F. M. Torres~\cite{lemos-silva_lotka-volterra_2023}, en el que analizan la inconsistencia de Euler explícito con el sistema continuo de Lotka-Volterra. uno de los argumentos que utilizan es que el método hace posible predecir densidades de población negativas, aunque el problema esté definido en $\mathbb{R}^2_+$. 
Para ello utiliza el ejemplo para los parámetros:
\begin{equation}
    \alpha = 1
    \label{eq:calc:alpha}
\end{equation}
\begin{equation}
    \beta = 0.1
    \label{eq:calc:beta}
\end{equation}
\begin{equation}
    \gamma = 0.075
    \label{eq:calc:gamma}
\end{equation}
\begin{equation}
    \delta = 0.75
    \label{eq:calc:delta}
\end{equation}
Tomamos además las condiciones iniciales siguientes
\begin{equation}
    x = 5
    \label{eq:calc:x}
\end{equation}
\begin{equation}
    y = 5
    \label{eq:calc:y}
\end{equation}
y el paso
\begin{equation}
    h = 0.03
    \label{eq:calc:h}
\end{equation}
Se evaluará en el tiempo $t=292.2$ , que con el paso $h$ anterior se corresponde con la iteración
\begin{equation}
    n = 9741
    \label{eq:calc:n}
\end{equation}

Ahora, aplicando el método de Euler al sistema, se obtienen las iteraciones
\[
\left \{
\begin{array}{rrrr}
    x_{i+1} & = &x_i + h(\alpha x_i - \beta x_i y_i) \\
    y_{i+1} & = &y_i + h(\gamma x_i y_i - \delta y_i)
\end{array}
\right .
\]
por lo que se obtiene
\begin{algorithm}[H]
\caption{Pseudocódigo de Euler explícito para el sistema de Lotka-Volterra}
\label{alg:calc:x}
\begin{algorithmic}
    \STATE $t \gets 0$
    \FOR{$i = 0$ \TO $n$}
        \STATE $t \gets h \cdot i$
        \STATE $dxdt \gets \alpha \cdot x - \beta \cdot x \cdot y$
        \STATE $dydt \gets \gamma \cdot x \cdot y - \delta \cdot y$
        \STATE $x \gets x + h \cdot dxdt$
        \STATE $y \gets y + h \cdot dydt$
    \ENDFOR
    \PRINT{t}
    \RETURN $x$
\end{algorithmic}
\end{algorithm}

Las poblaciones aproximadas en $t = 292.2$ obtenidas por el método de Euler ejecutado por ComPyLaTex son:

\begin{equation}
    x =
    -4.857146
    \label{eq:res:x}
\end{equation}
\begin{equation}
    y =
    504.588197
    \label{eq:res:y}
\end{equation}
En conclusión, esto demuestra que el método de Euler es inestable para el problema de Lotka-Volterra, ya que el problema está definido en $\mathbb{R}^2_+$.





\end{document}